\documentclass{article}
\usepackage[left=2cm, right=2cm, top=2cm, bottom=2cm]{geometry}
\usepackage{graphicx}
\usepackage{amsmath}
\usepackage{amssymb}
\usepackage{amsthm}
\usepackage{fancyhdr}
\usepackage{verbatim}
\usepackage{listings}
\usepackage{xcolor}
\usepackage{pdfpages}
\usepackage{cancel}
\usepackage{physics}
\usepackage{esint}
\usepackage{braket}

\lstset{
    language=C++,
    basicstyle=\ttfamily\footnotesize,
    keywordstyle=\color{blue},
    commentstyle=\color{green},
    stringstyle=\color{red},
    numbers=left,
    numberstyle=\tiny\color{gray},
    frame=single,
    breaklines=true,
    captionpos=b,
}


\title{MTH 446: Lecture \# 2 (catching up)}
\author{Cliff Sun}

\newtheorem{theorem}{Theorem}[section]
\newtheorem{lemma}[theorem]{Lemma}
\newtheorem{definition}[theorem]{Definition}
\newtheorem{conjecture}[theorem]{Conjecture}
\newtheorem{proposition}[theorem]{Proposition}
\newtheorem{corollary}[theorem]{Corollary}
\newtheorem{one minute paper}[theorem]{One Minute Paper}

\pagestyle{fancy}
\lhead{\textbf{\thepage}\ \ \nouppercase{\rightmark}}
\chead{MTH 446: Lecture \# 2 (catching up)}
\rhead{Cliff Sun}

\begin{document}

\maketitle

\section*{Lecture Span}
\begin{enumerate}
    \item Properties of $|z|$
\end{enumerate}

If $z = x + iy$, then $|z| = \sqrt{x^2 + y^2} = |\overline{z}|$. Some more properties:

\begin{enumerate}
    \item $\forall z \in \mathbb{C}$, $|z| \geq 0$ and $|z| = 0 \iff z = 0$
    \item $\forall z_1,z_2 \in \mathbb{C}$, $|z_1z_2| = |z_1||z_2|$
    \item (Triangle inequality), $\forall z_1z_2 \in \mathbb{C}$, $|z_1 + z_2| \leq |z_1| + |z_2|$
\end{enumerate}

\begin{proof}
    (Proof of 2). 
    \begin{align}
        |z_1z_2|^2 &= (z_1z_2)(\overline{z_1z_2}) \\
        &= (z_1z_2)(\overline{z_1}\overline{z_2}) \\
        &= (z_1 \overline{z_1})(z_2 \overline{z_2}) \\
        &= |z_1|^2|z_2|^2
    \end{align}
\end{proof}

\begin{proof}
    (Proof of 3)
    \begin{align}
        |z_1 + z_2|^2 &= (z_1 + z_2)(\overline{z_1 + z_2}) \\
        &= (z_1 + z_2)(\overline{z_1} + \overline{z_2}) \\
        &= |z_1|^2 + z_1\overline{z_2} + z_2\overline{z_1} + |z_2|^2\\
        &= |z_1|^2 + 2\Re(z_1z_2) + |z_2|^2\\
        &\leq |z_1|^2 + |z_1||z_2| + |z_2|^2\\
        &= (|z_1| + |z_2|)^2
    \end{align}
\end{proof}

\begin{corollary}
    $\forall z_1,z_2 \in \mathbb{C}$, $||z_1| - |z_2|| \leq |z_1 - z_2|$
\end{corollary}

\begin{proof}
    \begin{align}
        |z_1| &= |z_2 + (z_1 - z_2)| \\
        &\leq |z_2| + |z_1 - z_2| \\
        |z_1| - |z_2| \leq |z_1 - z_2|
    \end{align}
\end{proof}

You also can construct an alternative form of the equation of the circle centered at $z_0$ with radius $z$:

\begin{equation}
    |z - z_0| = R \iff |z|^2 - 2\Re(z \overline{z_0}) + |z_0|^2 = R^2
\end{equation}

Example: given $13 = 2^2 + 3^2$ and $74 = 5^2 + 7^2$, express $962 = 13 \cdot 74$ as the sum of two squares. 

\begin{proof}
    Let $z = 2 + 3i$ and $w = 5+ 7i$, then 
    \begin{align}
        |zw|^2 &= |z|^2|w|^2 \\
        &= 13 \cdot 74 = 962
    \end{align}
    Then calculate 
    \begin{equation}
        zw = -11 + 29i
    \end{equation}
    Thus, 
    \begin{equation}
        962 = 11^2 + 29^2
    \end{equation}
\end{proof}

\section*{Polar coordinates}

Let $z \in \mathbb{Z}$ and $z \neq 0$. Let $r = |z|$, then 

\begin{equation}
    \arg(z) = \{\theta + 2\pi n: n \in \mathbb{Z}\}
\end{equation}

This is called the argument of $z$. 

\end{document}